% ------------------------------------------------------------
% Section 3 -- Method
% ------------------------------------------------------------
\section{Method}
\label{sec:method}

% Fig. 1 -- architecture schematic (TikZ, spans both columns)
\begin{figure*}[t]
  \centering
  % ------------------------------------------------------------
% Fig. 1 -- Architecture schematic (TikZ)
% Input from sections/3_method.tex inside a figure* environment.
% ------------------------------------------------------------
\begin{tikzpicture}[
  font=\small,
  node distance=7mm and 9mm,
  box/.style={draw, rounded corners=2pt, align=center, minimum height=8.5mm, inner sep=4pt},
  trunk/.style={box, fill=blue!8},
  head/.style={box, fill=orange!12},
  loss/.style={box, fill=green!10, dashed},
  op/.style={draw, circle, inner sep=1.5pt},
  arr/.style={-{Stealth[length=2.2mm]}, thick},
  lbl/.style={font=\scriptsize, align=center}
]

% ---------- top row: trunk path ----------
\node[box] (x) {LDCT input\\$x$};
\node[trunk, right=of x] (B) {Denoising trunk $B$\\(RED-CNN / ResNet / any)\\ \scriptsize frozen interface: trunk unmodified};
\node[op, right=of B] (minus1) {$-$};
\node[box, right=of minus1] (nhat) {Removed noise\\$\hat{n} = x - B(x)$};

\draw[arr] (x) -- (B);
\draw[arr] (B) -- (minus1);
\draw[arr] (x.north) .. controls +(0,10mm) and +(0,10mm) .. (minus1.north)
  node[pos=0.5, lbl, above] {$x$};
\draw[arr] (minus1) -- (nhat);

% ---------- spectral head ----------
\node[head, below=11mm of nhat] (fft) {2-D FFT};
\node[head, left=of fft] (gain) {Radial gain $G_i(|f|)$\\ \scriptsize 32 log-gain knots \\ \scriptsize DC gain $=1$ (exact) \\ \scriptsize 2 near-DC knots frozen at $G{=}1$};
\node[head, left=of gain] (ifft) {inverse FFT};

\draw[arr] (nhat) -- (fft);
\draw[arr] (fft) -- (gain) node[pos=0.5, lbl, above] {$\times$};
\draw[arr] (gain) -- (ifft);

% ---------- conditioning encoder ----------
\node[head, below=10mm of gain] (enc) {Conditioning encoder ($\approx$5.3k params)\\ \scriptsize 3 strided convs $\to$ GAP $\to$ zero-init linear \\ \scriptsize $\Delta_i = \Delta_{\max}\tanh(\cdot)$, $|\Delta \log G| \le 0.25$};
\draw[arr] (enc) -- (gain) node[pos=0.5, lbl, right] {$\log G_i = \log G_{\mathrm{shared}} + \Delta_i$};
\draw[arr] (x.south) .. controls +(0,-24mm) and +(-14mm,0) .. (enc.west)
  node[pos=0.75, lbl, below] {$x$};

% ---------- output ----------
\node[op, left=13mm of ifft] (minus2) {$-$};
\node[box, left=of minus2, fill=gray!8] (xhat) {Output\\$\hat{x} = x - H(\hat{n};x)$};
\draw[arr] (ifft) -- (minus2) node[pos=0.5, lbl, above] {$H(\hat{n};x)$};
\draw[arr] (x.west) .. controls +(-8mm,0) and +(0,14mm) .. (minus2.north)
  node[pos=0.3, lbl, left] {$x$};
\draw[arr] (minus2) -- (xhat);

% ---------- losses ----------
\node[loss, below=10mm of enc, xshift=-12mm] (losses) {Training objective:\quad $L = L_{\mathrm{MSE}} + w_{\mathrm{nps}} L_{\mathrm{NPS}} + w_{\mathrm{hu}} L_{\mathrm{HUbin}} + w_c L_{\mathrm{center}}$\\[1pt] \scriptsize $L_{\mathrm{NPS}}$: radial log-NPS of $x-\hat{x}$ vs.\ real noise \quad $L_{\mathrm{HUbin}}$: per-tissue signed HU bias \quad $L_{\mathrm{center}}$: batch-centering of $\Delta_i$ (anti-degeneration)};

\end{tikzpicture}

  \caption{Overview of the proposed physics-faithful framework. A frozen-interface spectral head $H$ reshapes the removed-noise residual $\hat{n} = x - B(x)$ in the Fourier domain via a learnable radial gain curve $G(|f|)$ (32 knots; DC gain fixed at 1; two near-DC knots frozen at $G{=}1$). A lightweight conditioning encoder ($\approx$5.3k parameters) produces per-image, $\tanh$-bounded log-gain offsets $\Delta_i$, kept anatomy-discriminative by a batch-centering penalty. Training combines MSE with a radial log-NPS loss and a per-tissue HU-bin bias loss.}
  \label{fig:architecture}
\end{figure*}

\subsection{Trunk-agnostic residual interface}
\label{sec:interface}

Let $x$ be the LDCT input and $B$ any denoising trunk. We define the removed noise $\hat{n} = x - B(x)$ and produce the final output $\hat{x} = x - H(\hat{n};\, x)$, where $H$ is the spectral head. Because the interface requires only the pair $(x, B(x))$, the head attaches to any denoiser --- CNN, GAN generator, or transformer --- without modifying the trunk. All main results use a benchmark-faithful RED-CNN trunk as the representative testbed (Sec.~\ref{sec:related}).

\subsection{The DC-preserving radial spectral head}
\label{sec:head}

$H$ multiplies the 2-D Fourier transform of $\hat{n}$ by a learnable radial gain curve $G(|f|)$, parameterized by 32 log-gain knots interpolated over radial frequency bins. Four structural guarantees hold per image, by construction rather than by training:

\begin{enumerate}
  \item \textbf{Exact DC preservation.} The gain applies to non-zero frequencies only; the spatial mean (global HU level) of the residual is untouched. Provable, not learned.
  \item \textbf{Near-DC freezing.} The first two knots are pinned at $G = 1$, protecting large-area regional HU means. This constraint was introduced after diagnosing a $-7$\hu{} regional chest soft-tissue bias and faint concentric rings caused by a sharp gain step adjacent to DC (chest soft-tissue bias $-7.0 \rightarrow -3.5$\hu{} after freezing).
  \item \textbf{Radial parameterization.} The curve is 32 scalars --- isotropic, interpretable, directly plottable; the head cannot hallucinate anatomical structure.
  \item \textbf{Bounded adaptivity.} Per-image offsets pass through $\tanh$ with $|\Delta \log G| \le 0.25$; frozen bins receive zero offset.
\end{enumerate}

Parameter cost: 32 scalars (static) plus a ${\approx}5.3$k-parameter conditioning encoder (adaptive) --- $\le 10^{-3}$ of trunk parameters --- and one FFT pair per image at inference.

\subsection{Anatomy-adaptive conditioning}
\label{sec:conditioning}

A lightweight encoder (three strided convolutions $\rightarrow$ global average pooling $\rightarrow$ zero-initialized linear projection) maps the input image to per-bin log-gain offsets $\Delta_i$, yielding a per-image curve $G_i(|f|) = \exp(\log G_{\mathrm{shared}} + \Delta_i)$. Zero-initialization guarantees exact equivalence to the static head at initialization.

\subsection{A diagnosed failure mode and the centering constraint}
\label{sec:centering}

Trained naively, the adaptive offsets saturate at $-\Delta_{\max}$ for all images and all frequency bands: the adaptive path degenerates into a static global gain shift, which the trunk absorbs by denoising more aggressively --- image metrics improve (the selection metric is SSIM) while physics regresses (Sec.~\ref{sec:ablation}, config S4). Two general lessons follow: adaptive capacity serves the dominant training signal unless explicitly constrained, and any static spectral correction is trunk-absorbable under joint training --- only the between-image component of an adaptive correction is architecturally irreducible. We therefore add a quadratic \textbf{batch-centering penalty} $L_{\mathrm{center}}$ (weight 0.05) pulling the batch-mean adaptive offset to zero per frequency bin. Centering assigns the static component to the shared curve, reserves the adaptive path for between-image differences, and keeps $\tanh$ out of its zero-gradient region. With it, offsets remain unsaturated and become anatomy-discriminative (Sec.~\ref{sec:adaptivity}).

\subsection{Training objective}
\label{sec:objective}

\begin{equation}
L = L_{\mathrm{MSE}} + \wnps \cdot L_{\mathrm{NPS}} + \whu \cdot L_{\mathrm{HUbin}} + w_c \cdot L_{\mathrm{center}}.
\end{equation}

\paragraph{Radial NPS loss.}
$L_{\mathrm{NPS}}$ matches the radially averaged log noise-power spectrum of the removed noise $x - \hat{x}$ to that of the real noise $\mathrm{LDCT} - \mathrm{NDCT}$, penalizing both over-removal and the characteristic low-frequency under-removal hole left by appearance-trained trunks.

\paragraph{HU-bin bias loss.}
$L_{\mathrm{HUbin}}$ penalizes the mean signed error within fixed physical HU bins with boundaries $(-1024, -500, -200, 200, 600, 1900)$\hu{} --- air/lung, fat/low-density, soft tissue, dense tissue, bone --- directly targeting per-tissue calibration.

Pilot screening (17 configurations) established that $L_{\mathrm{NPS}}$ must always be paired with $L_{\mathrm{HUbin}}$ (alone, it redistributes HU error into soft tissue, bias ${\approx} -4$\hu{}) and that the useful region is $\wnps \in [0.005, 0.01]$; heavier weights ($\ge 0.03$) fail doubly (NPS worsens again, noise-magnitude ratio exceeds 1.1).
